% ============================================================================
% Chapter 36: How the Internet Works
% ============================================================================
\chapter{How the Internet Works}
\label{ch:internet}

\begin{objectives}
  \item Define the Internet and explain how it connects people
  \item Understand the roles of ISPs, routers, and servers
  \item Differentiate between the Internet and the World Wide Web
  \item Know what URLs, links, and web pages are
\end{objectives}

\bytetip[tip]{The Internet is the biggest network ever created --- it connects \textbf{billions} of devices across every country on Earth. Let's understand how this magic works!}

\section{What Is the Internet?}
\label{sec:what-internet}
\index{Internet}

\begin{theorybox}[Definition: The Internet]
The \textbf{Internet}\index{Internet} is a global network of interconnected computers that communicate using standard protocols (rules). It allows people to share information, communicate, and access services from anywhere in the world.
\end{theorybox}

% --- Figure 36.1: How You Connect to a Website ---
\begin{figure}[H]
\centering
\begin{tikzpicture}[node distance=1.2cm]
  \node[hwbox, fill=FunSky!10, minimum width=2cm] (you) {\faLaptop\\Your PC};
  \node[hwbox, fill=FunMint!10, draw=FunMint, minimum width=2cm, right=of you] (router) {\faWifi\\Router};
  \node[hwbox, fill=NoteOrange!10, draw=NoteOrange, minimum width=2cm, right=of router] (isp) {\faBuilding\\ISP};
  \node[hwbox, fill=FunPurple!10, draw=FunPurple, minimum width=2cm, right=of isp] (net) {\faGlobe\\Internet};
  \node[hwbox, fill=LabGreen!10, draw=LabGreen, minimum width=2cm, right=of net] (server) {\faServer\\Web Server};
  
  \draw[thickarrow] (you) -- (router);
  \draw[thickarrow] (router) -- (isp);
  \draw[thickarrow] (isp) -- (net);
  \draw[thickarrow] (net) -- (server);
  
  \node[below=0.5cm, font=\scriptsize\sffamily\itshape, text=NavyBlue, text width=14cm, align=center] at ($(isp)!0.5!(net)+(0,-0.3)$) {Your request travels through cables and satellites across the world to reach the server, and the response comes back the same way --- all in under a second!};
\end{tikzpicture}
\caption{How your computer connects to a website on the Internet}
\label{fig:internet-path}
\end{figure}

\section{Internet vs World Wide Web}
\label{sec:vs-www}
\index{World Wide Web}

\begin{itemize}[leftmargin=*, label={\textcolor{ModuleBlue}{\faAngleRight}}, itemsep=4pt]
  \item The \textbf{Internet}\index{Internet} is the \emph{physical network} --- the cables, routers, and satellites
  \item The \textbf{World Wide Web} (WWW)\index{WWW} is a \emph{service} that runs on the Internet --- the collection of websites you browse
  \item Think of the Internet as the \textbf{road system}, and the Web as the \textbf{shops and buildings} along those roads
\end{itemize}

\section{Understanding URLs}
\label{sec:urls}
\index{URL}

% --- Figure 36.2: URL Anatomy ---
\begin{figure}[H]
\centering
\begin{tikzpicture}
  \node[font=\large\ttfamily\bfseries, text=NavyBlue] at (6,2) {https://www.google.com/search};
  
  \draw[FunCoral, line width=1.5pt, ->] (1.2,2) -- (1.2,0.5);
  \node[below, font=\scriptsize\sffamily\bfseries, text=FunCoral, align=center] at (1.2,0.5) {\texttt{https://}\\Protocol\\(secure connection)};
  
  \draw[LabGreen, line width=1.5pt, ->] (4.5,2) -- (4.5,0.5);
  \node[below, font=\scriptsize\sffamily\bfseries, text=LabGreen, align=center] at (4.5,0.5) {\texttt{www.}\\World Wide\\Web prefix};
  
  \draw[FunPurple, line width=1.5pt, ->] (7.2,2) -- (7.2,0.5);
  \node[below, font=\scriptsize\sffamily\bfseries, text=FunPurple, align=center] at (7.2,0.5) {\texttt{google.com}\\Domain name\\(website address)};
  
  \draw[NoteOrange, line width=1.5pt, ->] (10.2,2) -- (10.2,0.5);
  \node[below, font=\scriptsize\sffamily\bfseries, text=NoteOrange, align=center] at (10.2,0.5) {\texttt{/search}\\Path\\(specific page)};
\end{tikzpicture}
\caption{Anatomy of a URL --- every part has a purpose}
\label{fig:url-anatomy}
\end{figure}

\bytetip[fun]{The first-ever website was created by \textbf{Tim Berners-Lee} on December 20, 1990. It described the World Wide Web project itself. The website is still online today at \texttt{info.cern.ch}!}

\begin{funfact}
There are currently over \textbf{2 billion websites} on the Internet, but only about 400 million are actively maintained. The rest are ``parked'' domain names or abandoned sites.
\end{funfact}

\begin{reviewqs}
  \item What is the Internet? Define it in your own words.
  \item What is the difference between the Internet and the World Wide Web?
  \item What does ISP stand for, and what does it do?
  \item What are the parts of a URL?
  \item Who invented the World Wide Web?
\end{reviewqs}

\begin{challenge}
Visit \textbf{5 different websites} and write down the full URL of each. For each URL, identify: the protocol, domain name, and path. Present your findings in a neat table.
\end{challenge}
